\chapter{Los cálculos que respaldan la placa}\label{cap:calculos}
\DeclareSIUnit\mil{mil}

En los capítulos anteriores aparecieron frases como ``esa pista aguanta de sobra'' o ``la caída es despreciable''.
Este capítulo muestra los números detrás de esas frases. Todo sale del archivo \path{hardware/pcb/CALCULOS.md}, que no se
escribe a mano: lo genera el script \path{tools/calculos.py} leyendo directamente la placa, \path{pcb.kicad_pcb}. Por eso
los largos, anchos, vías y separaciones que verás aquí son \emph{medidas de la placa real}, no suposiciones.

En este capítulo vas a ver:
\begin{itemize}
  \item de dónde salen los datos de entrada y qué tan confiable es cada uno;
  \item cuánta corriente lleva cada riel y por qué esos son los valores de peor caso;
  \item los tiempos de encendido y apagado del buck LMR36015 y el rizado de su bobina;
  \item cómo se calcula el ancho mínimo de una pista con IPC-2221, con un ejemplo completo;
  \item la caída de tensión en las pistas, la corriente por vía y la temperatura de cada pieza;
  \item la separación mínima entre pistas por tensión y la longitud crítica de los relojes;
  \item un resumen del análisis del op-amp TLV9152, que se explica completo en el capítulo~\ref{cap:analog}.
\end{itemize}

\begin{idea}[por qué un script y no una hoja de cálculo]
Si calculas a mano el largo de una pista y después mueves una vía, el cálculo queda viejo y nadie se entera.
\archivo{calculos.py} recorre las pistas de la placa con la API de KiCad (\texttt{pcbnew}), arma un grafo con cada red
y mide los caminos reales. Cada vez que la placa cambia, se vuelve a correr y \archivo{CALCULOS.md} queda al día. Los
números de este capítulo corresponden a la placa final (la misma que se exportó a Gerber).
\end{idea}

\section{Datos de entrada y qué tan confiables son}\label{sec:c10-datos}

Antes de calcular nada hay que saber con qué números se calcula. La tabla~\ref{tab:c10-datos} recoge los datos de
entrada con su etiqueta de origen. Recuerda el significado de las etiquetas:
\doc{} es un dato leído del datasheet (con la página), \copia{} viene de una fuente secundaria que no se contrastó
con el original y \estimado{} es un supuesto de diseño.

\begin{longtable}{@{}L{0.36\linewidth}L{0.32\linewidth}L{0.26\linewidth}@{}}
\caption{Datos de entrada de los cálculos (de \archivo{CALCULOS.md} §1).}\label{tab:c10-datos}\\
\toprule
\textbf{Dato} & \textbf{Valor} & \textbf{Origen}\\\midrule\endfirsthead
\toprule
\textbf{Dato} & \textbf{Valor} & \textbf{Origen}\\\midrule\endhead
Cobre externo (\capa{F.Cu}, \capa{B.Cu}) & \SI{35}{\micro\meter} (\SI{1}{\oz}) & apilado JLC04161H-7628\\
Cobre interno (\capa{In1.Cu}, \capa{In2.Cu}) & \SI{15.2}{\micro\meter} (\SI{0.5}{\oz}) & apilado JLC04161H-7628\\
Resistividad del cobre a \SI{20}{\degreeCelsius} & \SI{1.72e-8}{\ohm\meter} & \copia\\
Cobre en el barril de la vía & \SI{18}{\micro\meter} & \estimado{} (capacidad de JLC, sin verificar)\\
LMR36015, frecuencia (variante de \SI{1}{\mega\hertz}) & \num{0.85} / \num{1} / \SI{1.15}{\mega\hertz} (mín/típ/máx) & \doc{SNVSB49D p.7}\\
LMR36015, $t_\text{ON-MIN}$ / $t_\text{OFF-MIN}$ & \SI{83}{\nano\second} / \SI{73}{\nano\second} máx & \doc{SNVSB49D p.7}\\
LMR36015, $R_\text{DS-ON}$ alto / bajo & \SI{435}{\milli\ohm} / \SI{280}{\milli\ohm} máx & \doc{SNVSB49D p.7}\\
LMR36015, $R_{\theta JA}$ & \SI{72.5}{\degreeCelsius\per\watt} (JEDEC) & \doc{SNVSB49D p.6}\\
LMR36015 en FPWM & CCM sin carga, frecuencia fija & \doc{SNVSB49D p.16}\\
L1 (3V3) / L2 (5V) & \SI{6.8}{\micro\henry} / \SI{10}{\micro\henry} & \archivo{design.py}\\
PCM1808, $I_{CC}$ (VCC = \SI{5}{\volt}) & \SI{8.6}{\milli\ampere} típ, \SI{11}{\milli\ampere} máx & \doc{SLES177B p.6}\\
PCM1808, $I_{DD}$ (VDD = \SI{3.3}{\volt}) & \SI{5.9}{\milli\ampere} típ a \SI{48}{\kilo\hertz}; \SI{10.2}{\milli\ampere} típ a \SI{96}{\kilo\hertz} & \doc{SLES177B p.6}\\
TLV9152, $I_Q$ por amplificador & \SI{560}{\micro\ampere} típ; \SI{750}{\micro\ampere} máx (\SIrange{-40}{125}{\degreeCelsius}) & \doc{SBOS986E p.12}\\
MCP4725, $I_{DD}$ & \SI{400}{\micro\ampere} máx & \doc{DS22039D p.3}\\
Oscilador Y1 & $\le \SI{10}{\milli\ampere}$ & \estimado{} (sin datasheet a mano)\\
LT3092, $\theta_{JA}$ (TS8) & \SI{57}{\degreeCelsius\per\watt} & \copia\\
Corrientes IEPE & CH1 \SI{2.81}{} $\to$ \SI{3.98}{\milli\ampere}; CH2 \SI{1.51}{} $\to$ \SI{2.51}{\milli\ampere} (con JP) & \archivo{design.py}, cambio C9\\
D2 SMBJ26A & $V_\text{RWM}$ = \SI{26}{\volt}; $V_\text{BR}$ mín = \SI{28.9}{\volt} & \copia\\
IPC-2221 (ampacidad) & $I = k\,\Delta T^{0,44} A^{0,725}$; $k = \num{0.048}$ ext., \num{0.024} int. & \copia\\
\bottomrule
\end{longtable}

\begin{ojo}[lo que no está verificado]
Cuatro datos de la tabla no vienen del datasheet original: la resistividad del cobre (es un valor de manual), el
cobre del barril de la vía, la $\theta_{JA}$ del LT3092 y los valores del SMBJ26A. A esto se suman las constantes de
IPC-2221. Ninguno de ellos está cerca de un límite (verás márgenes de 10 veces o más), así que un error razonable en
esos datos no cambia ninguna conclusión. Igual quedan marcados para que no se presenten como hechos.
\end{ojo}

\subsection{Corriente de diseño de cada riel}

Para dimensionar pistas no sirve la corriente ``típica'': se usa una \textbf{cota superior}, es decir, un valor que
sabemos que la corriente real no va a pasar. \archivo{calculos.py} las arma sumando los máximos de cada carga:

\begin{calculo}[corrientes de peor caso]
\textbf{5V\_A} (lo alimentan el PCM1808 en su pin VCC y los dos TLV9152):
\begin{equation}
I_\text{5V\_A} = \underbrace{\SI{11}{\milli\ampere}}_{\text{PCM1808 } I_{CC}\text{ máx}}
 + \underbrace{2 \times 2 \times \SI{0.75}{\milli\ampere}}_{\text{2 chips} \times \text{2 amplificadores}}
 = \SI{14.0}{\milli\ampere}
\end{equation}
El riel \net{5V} lleva lo mismo, porque R33 solo alimenta a \net{5V_A}.

\textbf{3V3}:
\begin{equation}
I_\text{3V3} = \underbrace{\SI{10.2}{\milli\ampere}}_{\text{PCM1808 } I_{DD}\text{ a 96 kHz}}
 + \underbrace{2\times\SI{0.4}{\milli\ampere}}_{\text{2 MCP4725}}
 + \underbrace{2\times\frac{\SI{3.3}{\volt}}{\SI{4.7}{\kilo\ohm}}}_{\text{pull-ups I2C}}
 + \underbrace{\SI{10}{\milli\ampere}}_{\text{Y1 (estimado)}}
 = \SI{22.4}{\milli\ampere}
\end{equation}
Fíjate que se toma el PCM1808 a \SI{96}{\kilo\hertz} (\SI{10.2}{\milli\ampere}) aunque la placa muestrea a
\SI{48}{\kilo\hertz} (\SI{5.9}{\milli\ampere}): es una cota holgada a propósito. Las pull-ups se cuentan como si las
dos líneas I2C estuvieran en bajo todo el tiempo.

\textbf{24V\_PROT}: se toma \SI{50}{\milli\ampere}, el límite al que se configura la fuente de banco. Las dos fuentes
IEPE juntas piden como mucho $\SI{3.98}{} + \SI{2.51}{} = \SI{6.49}{\milli\ampere}$.
\end{calculo}

\section{El buck LMR36015: tiempos, rizado y disipación}\label{sec:c10-buck}

Un convertidor buck enciende un transistor durante un tiempo $t_\text{ON}$ y lo apaga durante $t_\text{OFF}$, a una
frecuencia fija $f_\text{SW}$. La fracción del periodo que pasa encendido es el \emph{ciclo útil} $D$. En un buck
ideal la salida es $V_\text{OUT} = D\cdot V_\text{IN}$. El capítulo~\ref{cap:alim} explica el circuito; aquí solo nos
interesa un límite del integrado.

\begin{idea}[el tiempo mínimo de encendido]
El LMR36015 no puede encender su transistor por menos de $t_\text{ON-MIN} = \SI{83}{\nano\second}$ (máximo, en el
peor chip) ni apagarlo por menos de $t_\text{OFF-MIN} = \SI{73}{\nano\second}$ \doc{SNVSB49D p.7}. Si para regular
necesitara un pulso más corto, el chip tendría que ``saltarse'' pulsos o bajar la frecuencia (en inglés,
\emph{frequency foldback}). Eso metería una frecuencia variable, posiblemente dentro de la banda de audio. Por eso hay
que comprobar que el pulso que se pide siempre es más largo que el mínimo.
\end{idea}

\begin{calculo}[$t_\text{ON}$ y $t_\text{OFF}$ en el peor caso]
Las fórmulas son
\begin{equation}\label{eq:c10-ton}
D = \frac{V_\text{OUT}}{V_\text{IN}}, \qquad t_\text{ON} = \frac{D}{f_\text{SW}}, \qquad t_\text{OFF} = \frac{1-D}{f_\text{SW}}.
\end{equation}
El pulso más corto aparece con la salida más baja (\SI{3.3}{\volt}), la entrada más alta y el oscilador más rápido.
La entrada más alta que puede llegar sin que conduzca el TVS es su ruptura mínima, \SI{28.9}{\volt} \copia, y el
oscilador rápido es \SI{1.15}{\mega\hertz} \doc{SNVSB49D p.7}:
\begin{align}
D &= \frac{\num{3.3}}{\num{28.9}} = \num{0.1142}\\
t_\text{ON} &= \frac{\num{0.1142}}{\SI{1.15e6}{\hertz}} = \SI{99}{\nano\second} \;>\; \SI{83}{\nano\second}\quad(\text{margen } \SI{16}{\nano\second})
\end{align}
Para $t_\text{OFF}$ el peor caso es al revés: salida alta, entrada baja. Con \SI{5}{\volt} y \SI{18}{\volt} a
\SI{1.15}{\mega\hertz}: $t_\text{OFF} = (1 - \num{0.278})/\SI{1.15e6}{\hertz} = \SI{628}{\nano\second}$, muy por
encima de \SI{73}{\nano\second}.
\end{calculo}

La tabla~\ref{tab:c10-ton} resume las esquinas. Las tensiones de entrada tienen un sentido: \SI{18}{\volt} es una
entrada baja de prueba, \SI{23.5}{\volt} es \SI{24}{\volt} menos la caída de D1, \SI{26}{\volt} es la tensión de
trabajo del TVS y \SI{28.9}{\volt} es su ruptura mínima.

\begin{table}[htbp]
\centering
\caption{Tiempos de encendido y apagado del LMR36015 en las esquinas de tolerancia (de \archivo{CALCULOS.md} §2).}\label{tab:c10-ton}
\small
\begin{tabular}{@{}cccrrrr@{}}
\toprule
$V_\text{OUT}$ & $V_\text{IN}$ & $f_\text{SW}$ & $t_\text{ON}$ & margen sobre \SI{83}{\nano\second} & $t_\text{OFF}$ & margen sobre \SI{73}{\nano\second}\\\midrule
\SI{3.3}{\volt} & \SI{18}{\volt} & \SI{1.15}{\mega\hertz} & \SI{159}{\nano\second} & $+\SI{76}{\nano\second}$ & \SI{710}{\nano\second} & $+\SI{637}{\nano\second}$\\
\SI{3.3}{\volt} & \SI{23.5}{\volt} & \SI{1.15}{\mega\hertz} & \SI{122}{\nano\second} & $+\SI{39}{\nano\second}$ & \SI{747}{\nano\second} & $+\SI{674}{\nano\second}$\\
\SI{3.3}{\volt} & \SI{26}{\volt} & \SI{1.15}{\mega\hertz} & \SI{110}{\nano\second} & $+\SI{27}{\nano\second}$ & \SI{759}{\nano\second} & $+\SI{686}{\nano\second}$\\
\SI{3.3}{\volt} & \SI{28.9}{\volt} & \SI{0.85}{\mega\hertz} & \SI{134}{\nano\second} & $+\SI{51}{\nano\second}$ & \SI{1042}{\nano\second} & $+\SI{969}{\nano\second}$\\
\textbf{\SI{3.3}{\volt}} & \textbf{\SI{28.9}{\volt}} & \textbf{\SI{1.15}{\mega\hertz}} & \textbf{\SI{99}{\nano\second}} & $\mathbf{+\SI{16}{\nano\second}}$ & \SI{770}{\nano\second} & $+\SI{697}{\nano\second}$\\
\SI{5.0}{\volt} & \SI{18}{\volt} & \SI{1.15}{\mega\hertz} & \SI{242}{\nano\second} & $+\SI{159}{\nano\second}$ & \textbf{\SI{628}{\nano\second}} & $\mathbf{+\SI{555}{\nano\second}}$\\
\SI{5.0}{\volt} & \SI{23.5}{\volt} & \SI{1.15}{\mega\hertz} & \SI{185}{\nano\second} & $+\SI{102}{\nano\second}$ & \SI{685}{\nano\second} & $+\SI{612}{\nano\second}$\\
\SI{5.0}{\volt} & \SI{28.9}{\volt} & \SI{1.15}{\mega\hertz} & \SI{150}{\nano\second} & $+\SI{67}{\nano\second}$ & \SI{719}{\nano\second} & $+\SI{646}{\nano\second}$\\
\bottomrule
\end{tabular}
\end{table}

\begin{ojo}[el margen de 16 ns es real pero justo]
Con \SI{24}{\volt} normales el margen es de \SIrange{39}{82}{\nano\second}. Los \SI{16}{\nano\second} solo aparecen si
la entrada sube a \SI{28.9}{\volt}, algo que no pasa con la fuente de banco a \SI{24}{\volt}. Además, las fórmulas
son ideales: con pérdidas, $D$ sube un poco y $t_\text{ON}$ también, lo que juega a favor.
\end{ojo}

\subsection{Rizado de la bobina en modo FPWM}

La variante usada (LMR36015\textbf{F}) trabaja en \emph{FPWM}: conmuta siempre a frecuencia fija, aunque no haya
carga \doc{SNVSB49D p.16}. La ventaja es que el ruido de conmutación queda en una frecuencia conocida, lejos del
audio. El costo es que por la bobina circula una corriente triangular aunque la carga consuma casi nada. Ese
triángulo tiene una amplitud pico a pico
\begin{equation}\label{eq:c10-ripple}
\Delta I_L = \frac{(V_\text{IN} - V_\text{OUT})\,V_\text{OUT}}{V_\text{IN}\, L\, f_\text{SW}}.
\end{equation}
El peor caso es con el oscilador lento (menos $f_\text{SW}$ en el denominador) y la entrada alta.

\begin{calculo}[rizado de L1, salida de 3,3 V]
Con $V_\text{IN} = \SI{28.9}{\volt}$, $L = \SI{6.8}{\micro\henry}$ y $f_\text{SW} = \SI{0.85}{\mega\hertz}$:
\begin{equation}
\Delta I_L = \frac{(\num{28.9} - \num{3.3})\times \num{3.3}}{\num{28.9}\times \SI{6.8e-6}{\henry}\times\SI{0.85e6}{\hertz}}
 = \SI{0.506}{\ampere}\ \text{p-p}.
\end{equation}
El pico de corriente en la bobina es la carga más la mitad del triángulo:
$I_\text{pico} = \SI{22.4}{\milli\ampere} + \SI{0.506}{\ampere}/2 = \SI{0.275}{\ampere}$.
El valor eficaz (RMS) es $\sqrt{I^2 + \Delta I_L^2/12} = \SI{0.148}{\ampere}$.
\end{calculo}

\begin{table}[htbp]
\centering
\caption{Rizado de las bobinas en FPWM con el oscilador a \SI{0.85}{\mega\hertz}.}\label{tab:c10-ripple}
\small
\begin{tabular}{@{}cccrrr@{}}
\toprule
Salida & $L$ & $V_\text{IN}$ & $\Delta I_L$ p-p & Pico & RMS\\\midrule
\SI{3.3}{\volt} & \SI{6.8}{\micro\henry} & \SI{23.5}{\volt} & \SI{0.491}{\ampere} & \SI{0.268}{\ampere} & \SI{0.143}{\ampere}\\
\SI{3.3}{\volt} & \SI{6.8}{\micro\henry} & \SI{28.9}{\volt} & \SI{0.506}{\ampere} & \SI{0.275}{\ampere} & \SI{0.148}{\ampere}\\
\SI{5.0}{\volt} & \SI{10}{\micro\henry} & \SI{23.5}{\volt} & \SI{0.463}{\ampere} & \SI{0.246}{\ampere} & \SI{0.134}{\ampere}\\
\SI{5.0}{\volt} & \SI{10}{\micro\henry} & \SI{28.9}{\volt} & \SI{0.486}{\ampere} & \SI{0.257}{\ampere} & \SI{0.141}{\ampere}\\
\bottomrule
\end{tabular}
\end{table}

\begin{idea}[la sorpresa del buck sin carga]
La carga de \net{3V3} es de unos \SI{22}{\milli\ampere}, pero el tramo SW--bobina--condensador de salida lleva picos
de \SI{0.28}{\ampere}. Por eso esos tramos se dimensionan con \SI{0.28}{\ampere} y no con la corriente de la carga.
Es un detalle que se pasa por alto fácil si solo miras el consumo.
\end{idea}

\subsection{Disipación del LMR36015}

El datasheet no da la corriente de entrada en FPWM sin carga, así que la disipación se acota con partes conocidas y
supuestos \estimado:
\begin{itemize}
  \item \textbf{Conducción}: $P = I_\text{RMS}^2\,[D\,R_\text{DS-ON,alto} + (1-D)\,R_\text{DS-ON,bajo}]$ con las
  resistencias máximas: $\num{0.148}^2 \times [\num{0.114}\times\num{0.435} + \num{0.886}\times\num{0.280}] \approx \SI{6.5}{\milli\watt}$.
  \item \textbf{Conmutación}, suponiendo flancos de \SI{10}{\nano\second}:
  $P \approx \tfrac12 V_\text{IN}\,(\Delta I_L/2)\,t_\text{flanco}\,f_\text{SW} \approx \SI{42}{\milli\watt}$.
  \item El regulador interno y la carga de las compuertas no están publicados.
\end{itemize}
Con una cota de \SI{0.15}{\watt} por integrado, la subida de temperatura con la $R_{\theta JA}$ JEDEC es
$\Delta T = \SI{0.15}{\watt}\times\SI{72.5}{\degreeCelsius\per\watt} \approx \SI{11}{\degreeCelsius}$. La placa de 4
capas con el plano de tierra entero disipa mejor que la placa de prueba JEDEC, así que el valor real debería quedar
por debajo. \textbf{Esto se confirma en el banco} midiendo la corriente de entrada a \SI{24}{\volt} sin carga
(capítulo~\ref{cap:banco}).

\section{Ancho de pista por corriente (IPC-2221)}\label{sec:c10-ancho}

Una pista es una resistencia: al pasarle corriente se calienta. La norma IPC-2221 da una fórmula empírica (sale de
ajustar curvas medidas) que relaciona la corriente $I$ con la sección de cobre $A$ y el aumento de temperatura
$\Delta T$ que aceptas:
\begin{equation}\label{eq:c10-ipc}
I = k\,\Delta T^{\,0,44}\,A^{0,725},\qquad k = \num{0.048}\ \text{(capa externa)},\quad k = \num{0.024}\ \text{(capa interna)},
\end{equation}
con $I$ en amperios, $\Delta T$ en \si{\degreeCelsius} y $A$ en \si{\mil\squared} (\SI{1}{\mil} = \SI{0.0254}{\milli\meter})
\copia. Las capas internas tienen la mitad de $k$ porque están encerradas en el FR-4 y se enfrían peor. En este
proyecto se aceptó $\Delta T = \SI{10}{\degreeCelsius}$.

\begin{calculo}[ancho mínimo para el nodo SW (0,275 A, capa externa)]
\textbf{Paso 1}, despejar el área:
\begin{equation}
A = \left(\frac{I}{k\,\Delta T^{0,44}}\right)^{1/0,725}.
\end{equation}
\textbf{Paso 2}, el factor de temperatura: $\Delta T^{0,44} = 10^{0,44} = \num{2.754}$, y
$k\,\Delta T^{0,44} = \num{0.048}\times\num{2.754} = \num{0.1322}$.

\textbf{Paso 3}, el área:
\begin{equation}
A = \left(\frac{\num{0.275}}{\num{0.1322}}\right)^{1/0,725} = \num{2.080}^{\,\num{1.379}} = \SI{2.75}{\mil\squared}.
\end{equation}
\textbf{Paso 4}, pasar a ancho dividiendo por el espesor. \SI{35}{\micro\meter} son \SI{1.378}{\mil}:
\begin{equation}
w = \frac{\SI{2.75}{\mil\squared}}{\SI{1.378}{\mil}} = \SI{1.99}{\mil} = \SI{0.051}{\milli\meter}.
\end{equation}
La pista más angosta del nodo SW mide \SI{0.25}{\milli\meter} \medido, casi cinco veces más.

\textbf{Al revés}, ¿cuánto aguanta una pista de \SI{0.25}{\milli\meter} en capa externa?
$A = \num{0.25}\times\SI{0.035}{\milli\meter\squared} = \SI{13.6}{\mil\squared}$ y
$I = \num{0.1322}\times\num{13.6}^{0,725} \approx \SI{0.88}{\ampere}$.
Una pista de \SI{0.4}{\milli\meter} en \capa{In2.Cu} (\SI{15.2}{\micro\meter}, $k=\num{0.024}$) aguanta
unos \SI{0.34}{\ampere}.
\end{calculo}

\begin{longtable}{@{}L{0.30\linewidth}rcrL{0.28\linewidth}@{}}
\caption{Ancho mínimo IPC-2221 frente al ancho usado en la placa (de \archivo{CALCULOS.md} §3).}\label{tab:c10-anchos}\\
\toprule
\textbf{Tramo} & \textbf{Corriente} & \textbf{Capa} & \textbf{Mínimo IPC} & \textbf{Ancho mínimo usado} \medido\\\midrule\endfirsthead
\toprule
\textbf{Tramo} & \textbf{Corriente} & \textbf{Capa} & \textbf{Mínimo IPC} & \textbf{Ancho mínimo usado} \medido\\\midrule\endhead
\net{SW_3}/\net{SW_5}, bobina, COUT & \SI{275}{\milli\ampere} & \capa{F.Cu} & \SI{0.051}{\milli\meter} & \SI{0.25}{\milli\meter} (máx. \num{0.60}; \SI{12}{\milli\meter} de pista)\\
\net{24V_PROT} (bus, VIN) & \SI{50}{\milli\ampere} & \capa{F.Cu} & \SI{0.005}{\milli\meter} & \SI{0.25}{\milli\meter} (máx. \num{1.00}; \SI{98}{\milli\meter})\\
\net{3V3} & \SI{22}{\milli\ampere} & \capa{B.Cu} & \SI{0.002}{\milli\meter} & \SI{0.30}{\milli\meter} (\SI{4}{\milli\meter})\\
\net{3V3} & \SI{22}{\milli\ampere} & \capa{F.Cu} & \SI{0.002}{\milli\meter} & \SI{0.25}{\milli\meter} (\SI{52}{\milli\meter})\\
\net{3V3} & \SI{22}{\milli\ampere} & \capa{In2.Cu} & \SI{0.010}{\milli\meter} & \SI{0.40}{\milli\meter} (\SI{117}{\milli\meter})\\
\net{5V} (salida a R33) & \SI{14}{\milli\ampere} & \capa{F.Cu} & \SI{0.001}{\milli\meter} & \SI{0.25}{\milli\meter} (\SI{29}{\milli\meter})\\
\net{5V_A} & \SI{14}{\milli\ampere} & \capa{F.Cu} & \SI{0.001}{\milli\meter} & \SI{0.30}{\milli\meter} (\SI{26}{\milli\meter})\\
\net{5V_A} & \SI{14}{\milli\ampere} & \capa{In2.Cu} & \SI{0.005}{\milli\meter} & \SI{0.40}{\milli\meter} (\SI{53}{\milli\meter})\\
\net{IEPE_SUP}, \net{IEPE_SUP_SV} & \SI{4}{\milli\ampere} & \capa{F.Cu} & $\approx 0$ & \SI{0.30}{\milli\meter} (\SI{23}{\milli\meter})\\
\net{GND} (pistas) & \SI{50}{\milli\ampere} & \capa{F.Cu} & \SI{0.005}{\milli\meter} & \SI{0.25}{\milli\meter} (\SI{110}{\milli\meter})\\
\bottomrule
\end{longtable}

\begin{decision}[anchos de pista]
\textbf{Qué se decidió:} las pistas de potencia van de \SIrange{0.25}{0.6}{\milli\meter} en capas externas y a
\SI{0.4}{\milli\meter} en \capa{In2.Cu}.\par
\textbf{Alternativas:} usar el ancho por defecto de la clase Potencia (\SI{0.5}{\milli\meter}) en todas partes, o
pistas gruesas ``por si acaso''.\par
\textbf{Por qué:} con corrientes de \SIrange{5}{20}{\milli\ampere} (y \SI{0.28}{\ampere} de pico en SW), el ancho lo
fija lo que se puede fabricar y la regla mínima de la clase (\SI{0.25}{\milli\meter}), no la corriente. El margen es
de diez veces o más. Pistas más anchas no mejoran nada y quitan espacio en una placa apretada. Las salidas de pin van a
\SI{0.25}{\milli\meter} porque el pad del integrado tiene ese ancho.\par
\textbf{Fuente:} \archivo{CALCULOS.md} §3; IPC-2221 \copia.
\end{decision}

\section{Caída de tensión en las pistas}\label{sec:c10-caida}

Que una pista no se caliente no significa que no baje tensión. La resistencia de un tramo de cobre es
\begin{equation}\label{eq:c10-r}
R = \rho\,\frac{\ell}{w\,t},
\end{equation}
con $\rho$ la resistividad, $\ell$ el largo, $w$ el ancho y $t$ el espesor. Algunos valores por milímetro de largo:

\begin{itemize}
  \item pista de \SI{0.25}{\milli\meter} en capa externa (\SI{35}{\micro\meter}): \SI{1.97}{\milli\ohm} por milímetro;
  \item pista de \SI{0.3}{\milli\meter} en capa externa: \SI{1.64}{\milli\ohm} por milímetro;
  \item pista de \SI{0.4}{\milli\meter} en capa interna (\SI{15.2}{\micro\meter}): \SI{2.83}{\milli\ohm} por milímetro;
  \item una vía (barril de \SI{0.3}{\milli\meter}, \SI{18}{\micro\meter}, \SI{1.6}{\milli\meter} de alto): \SI{1.6}{\milli\ohm}.
\end{itemize}

Fíjate en el detalle de la capa interna: aunque la pista es más ancha, el cobre es menos de la mitad de grueso, así que
tiene más resistencia por milímetro.

\archivo{calculos.py} arma un grafo con todas las pistas y vías de cada red y busca el camino de menor resistencia
entre la fuente y la carga más lejana. Después multiplica por la corriente \emph{total} del riel, como si toda pasara
por el camino más largo (otra cota de peor caso).

\begin{table}[htbp]
\centering
\caption{Caída de tensión de la fuente a la carga más lejana \medido.}\label{tab:c10-caida}
\small
\begin{tabular}{@{}llllrrr@{}}
\toprule
Riel & Desde & Hasta & Largo & $R$ & $I$ & Caída\\\midrule
\net{3V3} & zona VOUT del buck & \net{U4.4} & \SI{82.9}{\milli\meter} & \SI{199}{\milli\ohm} & \SI{22.4}{\milli\ampere} & \SI{4.47}{\milli\volt}\\
\net{5V_A} & \net{R33.2} & \net{U9.8} & \SI{54.1}{\milli\meter} & \SI{125}{\milli\ohm} & \SI{14.0}{\milli\ampere} & \SI{1.75}{\milli\volt}\\
\bottomrule
\end{tabular}
\end{table}

\begin{calculo}[¿es mucho 4,47 mV?]
La salida del buck tiene una tolerancia de $\pm\SI{1.5}{\percent}$ en FPWM \doc{SNVSB49D p.8}. En \SI{3.3}{\volt}
eso es $\pm\SI{49.5}{\milli\volt}$. La caída en el cobre es menos de la décima parte de esa tolerancia. Comparado con
las resistencias que están a propósito en el camino, es todavía menor:
\begin{itemize}
  \item R33 (\SI{10}{\ohm}) baja \net{5V_A} en $\SI{14}{\milli\ampere}\times\SI{10}{\ohm} = \SI{140}{\milli\volt}$ y
  disipa $(\SI{14}{\milli\ampere})^2\times\SI{10}{\ohm} = \SI{2.0}{\milli\watt}$;
  \item R9 y R46 (\SI{100}{\ohm}) bajan cada \net{3V3_DAC} en \SI{40}{\milli\volt} con el MCP4725 al máximo.
\end{itemize}
Conclusión: el ancho de las pistas de alimentación no afecta la tensión que llega a los integrados.
\end{calculo}

\section{Corriente por vía}\label{sec:c10-vias}

Una vía es un tubito de cobre que atraviesa la placa. Su sección es la de un anillo delgado:
$A = \pi\,d\,t = \pi\times\SI{0.3}{\milli\meter}\times\SI{0.018}{\milli\meter} = \SI{0.0170}{\milli\meter\squared}$
(\SI{26.3}{\mil\squared}). Aplicando~\eqref{eq:c10-ipc} con $k=\num{0.048}$ (práctica común para el barril) da
\SI{1.41}{\ampere} por vía; con $k=\num{0.024}$, \SI{0.71}{\ampere}. En los dos casos, una sola vía basta para
cualquier riel de la placa.

\begin{table}[htbp]
\centering
\caption{Vías por red frente a las necesarias por corriente.}\label{tab:c10-vias}
\small
\begin{tabular}{@{}lrrr@{}}
\toprule
Red & Vías en la placa \medido & Corriente del riel & Vías necesarias\\\midrule
\net{3V3} & 11 & \SI{22.4}{\milli\ampere} & 1\\
\net{5V_A} & 5 & \SI{14.0}{\milli\ampere} & 1\\
\net{5V} & 2 & \SI{14.0}{\milli\ampere} & 1\\
\net{24V_PROT} & 0 & \SI{50}{\milli\ampere} & 1\\
\net{GND} & 299 & \SI{50}{\milli\ampere} & 1\\
\net{I2C_SCL} / \net{I2C_SDA} & 3 / 2 & \SI{1}{\milli\ampere} & 1\\
\net{DC_MON} / \net{DC_MON2} & 1 / 2 & \SI{0.1}{\milli\ampere} & 1\\
\bottomrule
\end{tabular}
\end{table}

Entonces, ¿por qué hay 299 vías de tierra si una alcanza? Porque no están ahí por corriente: unen el relleno de
tierra de \capa{F.Cu} con los planos de \capa{In1.Cu} y \capa{B.Cu} (cosido), dan a cada condensador de desacoplo un
camino corto al plano y evitan islas de cobre. Eso se explica en los capítulos~\ref{cap:capas} y~\ref{cap:ruteo}.
\net{24V_PROT} tiene 0 vías porque toda su distribución va en \capa{F.Cu}.

\section{Disipación y temperatura de cada pieza}\label{sec:c10-termico}

Para cada pieza que disipa algo se busca su peor caso, se calcula la potencia y se multiplica por su resistencia
térmica $\theta_{JA}$ (cuántos grados sube por cada vatio):
\begin{equation}
\Delta T = P\,\theta_{JA}.
\end{equation}

\begin{calculo}[el LT3092 con el micrófono en cortocircuito]
El peor caso de la fuente de corriente IEPE es que el micrófono o el cable estén en corto: la salida queda en
\SI{0}{\volt} y toda la tensión de alimentación cae en el LT3092. Con la corriente mayor (CH1 con JP1,
\SI{3.98}{\milli\ampere}) y \SI{23.3}{\volt} entre entrada y salida:
\begin{equation}
P = \SI{3.98}{\milli\ampere}\times\SI{23.3}{\volt} = \SI{93}{\milli\watt},\qquad
\Delta T = \SI{0.093}{\watt}\times\SI{57}{\degreeCelsius\per\watt} = \SI{5.3}{\degreeCelsius}.
\end{equation}
Un corto en el micrófono no daña la placa.
\end{calculo}

\begin{table}[htbp]
\centering
\caption{Disipación y subida de temperatura en el peor caso (de \archivo{CALCULOS.md} §6).}\label{tab:c10-termico}
\small
\begin{tabular}{@{}lL{0.27\linewidth}rrL{0.22\linewidth}@{}}
\toprule
Pieza & Peor caso & Potencia & $\Delta T$ & Nota\\\midrule
LT3092 (U1, U8) & IEPE en corto, \SI{3.98}{\milli\ampere} $\times$ \SI{23.3}{\volt} & \SI{93}{\milli\watt} & \SI{5.3}{\degreeCelsius} & $\theta_{JA}$ \copia\\
LMR36015 (U6, U7) & FPWM sin carga, \SI{28.9}{\volt} & $\le\SI{150}{\milli\watt}$ & $\le\SI{11}{\degreeCelsius}$ & \estimado\newline \medir\\
D1 SS34 & \SI{50}{\milli\ampere} $\times$ \SI{0.4}{\volt} & \SI{20}{\milli\watt} & $<\SI{1}{\degreeCelsius}$ & $V_F$ \estimado\\
R33 \SI{10}{\ohm} & \SI{14}{\milli\ampere} & \SI{2}{\milli\watt} & despreciable & 0603 de \SI{100}{\milli\watt}\\
D2 SMBJ26A & en reposo (solo fuga) & despreciable & --- & conduce solo en transitorios\\
PCM1808 & \SI{48}{\kilo\hertz} & \SIrange{62}{81}{\milli\watt} & --- & \doc{SLES177B p.6}\\
\bottomrule
\end{tabular}
\end{table}

Ninguna pieza se acerca a una temperatura preocupante. La de mayor incertidumbre es el LMR36015, y por eso es la
primera medición de la lista de banco.

\section{Separación entre pistas por tensión}\label{sec:c10-separacion}

Dos conductores con tensiones distintas necesitan aire (o FR-4) entre ellos para que no salte un arco ni haya fugas.
IPC-2221B, tabla 6-1, da la separación mínima según la tensión y el tipo de conductor \copia:

\begin{itemize}
  \item hasta \SI{30}{\volt}, capa externa sin recubrimiento, hasta \SI{3050}{\meter} de altura (categoría B2):
  \SI{0.1}{\milli\meter}. Manizales está a unos \SI{2150}{\meter}, dentro de esa categoría;
  \item hasta \SI{30}{\volt}, capa interna (B1): \SI{0.05}{\milli\meter};
  \item de \SIrange{31}{50}{\volt}, capa externa: \SI{0.6}{\milli\meter}.
\end{itemize}

¿En qué fila cae la placa? La tensión más alta es la de \SI{24}{\volt}, y el TVS D2 limita la tensión continua a su
$V_\text{RWM}$ de \SI{26}{\volt}. Así que las redes de \SI{24}{\volt} van en la fila de hasta \SI{30}{\volt}.

\begin{calculo}[separación mínima medida]
\archivo{calculos.py} toma todas las redes que pueden llegar a \SIrange{24}{29}{\volt}: \net{24V_PROT},
\net{IEPE_SUP}, \net{IEPE_SUP_SV}, \net{SW_3}, \net{SW_5}, \net{IEPE}, \net{IEPE_SV}, el ánodo de D1 y los pines 2 de
J1 y J7. Mide la distancia de su cobre a cualquier otra red, en todas las capas, con la geometría real de pads,
pistas y zonas. Resultado \medido:
\[
\text{separación mínima} = \SI{0.200}{\milli\meter}\ (\net{24V_PROT} \text{ contra } \net{GND}, \text{ en } \capa{F.Cu}),
\]
el doble de los \SI{0.1}{\milli\meter} que pide la norma para capa externa.
\end{calculo}

\section{Relojes: longitud crítica e impedancia}\label{sec:c10-relojes}

Cuando una pista es larga comparada con lo rápido que cambia la señal, deja de comportarse como un cable ideal y se
vuelve una \emph{línea de transmisión}: aparecen reflexiones y ``timbre'' (oscilaciones en los flancos). La regla
práctica dice que eso importa cuando el viaje de ida y vuelta por la pista dura más que el flanco. De ahí sale la
\emph{longitud crítica}
\begin{equation}\label{eq:c10-lc}
L_c = \frac{t_r}{2\,t_{pd}},
\end{equation}
con $t_r$ el tiempo de subida del flanco y $t_{pd}$ el retardo por milímetro de pista. En microstrip sobre FR-4
($\varepsilon_r \approx \num{4.4}$), $t_{pd} \approx \SI{5.6}{\pico\second\per\milli\meter}$ \copia. Con un flanco de
\SI{1}{\nano\second} (supuesto conservador, \estimado):
\[
L_c = \frac{\SI{1000}{\pico\second}}{2\times\SI{5.6}{\pico\second\per\milli\meter}} \approx \SI{89}{\milli\meter}.
\]

\begin{table}[htbp]
\centering
\caption{Largo de las líneas de reloj y datos frente a la longitud crítica \medido.}\label{tab:c10-relojes}
\small
\begin{tabular}{@{}lL{0.38\linewidth}rr@{}}
\toprule
Red & Qué es & Largo & \% de $L_c$\\\midrule
\net{Net-(Y1-OUT)} & salida del oscilador & \SI{4.2}{\milli\meter} & \SI{5}{\percent}\\
\net{Net-(U4-SCKI)} & reloj maestro al PCM1808 & \SI{5.4}{\milli\meter} & \SI{6}{\percent}\\
\net{Net-(U4-BCK)} & reloj de bit (lado ADC) & \SI{1.9}{\milli\meter} & \SI{2}{\percent}\\
\net{SAI_SCK} & reloj de bit (lado NUCLEO) & \SI{13.5}{\milli\meter} & \SI{15}{\percent}\\
\net{Net-(U4-LRCK)} & reloj de trama (lado ADC) & \SI{6.1}{\milli\meter} & \SI{7}{\percent}\\
\net{SAI_FS} & reloj de trama (lado NUCLEO) & \SI{25.7}{\milli\meter} & \SI{29}{\percent}\\
\net{Net-(U4-DOUT)} & datos (lado ADC) & \SI{4.2}{\milli\meter} & \SI{5}{\percent}\\
\net{SAI_SD} & datos (lado NUCLEO) & \SI{15.5}{\milli\meter} & \SI{17}{\percent}\\
\bottomrule
\end{tabular}
\end{table}

La más larga (\net{SAI_FS}) es menos de un tercio de $L_c$. Además, cada línea tiene una resistencia en serie (R21--R24)
junto a su fuente, que suaviza los flancos y amortigua cualquier rebote.

\begin{calculo}[impedancia de una pista de 0,2 mm]
Para un microstrip (pista sobre un plano de tierra) la aproximación de IPC-2141 \copia{} es
\begin{equation}
Z_0 = \frac{87}{\sqrt{\varepsilon_r + \num{1.41}}}\,\ln\!\left(\frac{\num{5.98}\,h}{\num{0.8}\,w + t}\right),
\end{equation}
con $h = \SI{0.2104}{\milli\meter}$ (prepreg 7628 entre \capa{F.Cu} e \capa{In1.Cu}), $w = \SI{0.2}{\milli\meter}$ y
$t = \SI{0.035}{\milli\meter}$:
\[
Z_0 = \frac{87}{\sqrt{\num{5.81}}}\,\ln\!\left(\frac{\num{1.258}}{\num{0.195}}\right) = \num{36.1}\times\num{1.864} \approx \SI{67}{\ohm}.
\]
Como todas las líneas son mucho más cortas que $L_c$, este valor no hace falta controlarlo: por eso el pedido a JLC no
lleva control de impedancia. En la placa no hay USB ni pares diferenciales.
\end{calculo}

\section{El op-amp TLV9152 (resumen)}\label{sec:c10-tlv}

La sección 9 de \archivo{CALCULOS.md} trae la salida de \archivo{tools/analisis_opamp.py}. El análisis completo, con
la explicación del circuito, está en el capítulo~\ref{cap:analog}. Los resultados clave son:

\begin{itemize}
  \item \textbf{Sumador U3A/U9A}: sin el condensador de realimentación el margen de fase sería de unos
  \SI{8}{\degree} (casi oscilando, con \SI{17}{\decibel} de pico). Con CF = \SI{10}{\pico\farad} (C56, C57) queda en
  unos \SI{70}{\degree}, sin pico, y la señal tiene su $-\SI{3}{\decibel}$ en \SI{166}{\kilo\hertz}, muy por encima
  del audio.
  \item \textbf{Seguidor U3B/U9B}: su carga son unos \SI{5}{\pico\farad} y el datasheet \doc{SBOS986E p.12}
  admite \SI{1000}{\pico\farad}. Es estable.
  \item \textbf{Región de transición de la entrada}: para quedar fuera de ella se recomienda (pendiente de aprobación)
  que el firmware lleve \net{REST_OUT} a unos \SI{1.5}{\volt}. Así \net{VSHIFT} queda en \SIrange{2.1}{2.6}{\volt}.
  \item \textbf{Sin verificar}: simulación con el modelo SPICE de TI y medición en banco.
\end{itemize}

\section*{Resumen}
\begin{itemize}
  \item Los cálculos salen de un script que mide la placa real; si la placa cambia, se regeneran.
  \item Las corrientes son cotas de peor caso: \SI{14}{\milli\ampere} en \net{5V_A}, \SI{22.4}{\milli\ampere} en
  \net{3V3} y \SI{50}{\milli\ampere} en \net{24V_PROT}.
  \item El buck nunca pide un pulso más corto que su mínimo: el peor caso deja \SI{16}{\nano\second} de margen, y con
  \SI{24}{\volt} normales hay \SIrange{39}{82}{\nano\second}.
  \item En FPWM la bobina lleva \SI{0.28}{\ampere} de pico aunque la carga sea de miliamperios.
  \item Todas las pistas tienen diez veces o más del ancho que pide IPC-2221; las caídas son de pocos milivoltios y una
  vía aguanta más de un amperio.
  \item Ninguna pieza sube más de unos \SI{11}{\degreeCelsius}; el LMR36015 es el dato a medir.
  \item La separación mínima de \SI{24}{\volt} es \SI{0.200}{\milli\meter}, el doble de lo que pide IPC-2221B.
  \item Los relojes son cortos frente a su longitud crítica, así que no hace falta controlar impedancia.
\end{itemize}
